% Options for packages loaded elsewhere
\PassOptionsToPackage{unicode}{hyperref}
\PassOptionsToPackage{hyphens}{url}
\documentclass[
]{article}
\usepackage{xcolor}
\usepackage[margin=1in]{geometry}
\usepackage{amsmath,amssymb}
\setcounter{secnumdepth}{-\maxdimen} % remove section numbering
\usepackage{iftex}
\ifPDFTeX
  \usepackage[T1]{fontenc}
  \usepackage[utf8]{inputenc}
  \usepackage{textcomp} % provide euro and other symbols
\else % if luatex or xetex
  \usepackage{unicode-math} % this also loads fontspec
  \defaultfontfeatures{Scale=MatchLowercase}
  \defaultfontfeatures[\rmfamily]{Ligatures=TeX,Scale=1}
\fi
\usepackage{lmodern}
\ifPDFTeX\else
  % xetex/luatex font selection
\fi
% Use upquote if available, for straight quotes in verbatim environments
\IfFileExists{upquote.sty}{\usepackage{upquote}}{}
\IfFileExists{microtype.sty}{% use microtype if available
  \usepackage[]{microtype}
  \UseMicrotypeSet[protrusion]{basicmath} % disable protrusion for tt fonts
}{}
\makeatletter
\@ifundefined{KOMAClassName}{% if non-KOMA class
  \IfFileExists{parskip.sty}{%
    \usepackage{parskip}
  }{% else
    \setlength{\parindent}{0pt}
    \setlength{\parskip}{6pt plus 2pt minus 1pt}}
}{% if KOMA class
  \KOMAoptions{parskip=half}}
\makeatother
\usepackage{color}
\usepackage{fancyvrb}
\newcommand{\VerbBar}{|}
\newcommand{\VERB}{\Verb[commandchars=\\\{\}]}
\DefineVerbatimEnvironment{Highlighting}{Verbatim}{commandchars=\\\{\}}
% Add ',fontsize=\small' for more characters per line
\usepackage{framed}
\definecolor{shadecolor}{RGB}{248,248,248}
\newenvironment{Shaded}{\begin{snugshade}}{\end{snugshade}}
\newcommand{\AlertTok}[1]{\textcolor[rgb]{0.94,0.16,0.16}{#1}}
\newcommand{\AnnotationTok}[1]{\textcolor[rgb]{0.56,0.35,0.01}{\textbf{\textit{#1}}}}
\newcommand{\AttributeTok}[1]{\textcolor[rgb]{0.13,0.29,0.53}{#1}}
\newcommand{\BaseNTok}[1]{\textcolor[rgb]{0.00,0.00,0.81}{#1}}
\newcommand{\BuiltInTok}[1]{#1}
\newcommand{\CharTok}[1]{\textcolor[rgb]{0.31,0.60,0.02}{#1}}
\newcommand{\CommentTok}[1]{\textcolor[rgb]{0.56,0.35,0.01}{\textit{#1}}}
\newcommand{\CommentVarTok}[1]{\textcolor[rgb]{0.56,0.35,0.01}{\textbf{\textit{#1}}}}
\newcommand{\ConstantTok}[1]{\textcolor[rgb]{0.56,0.35,0.01}{#1}}
\newcommand{\ControlFlowTok}[1]{\textcolor[rgb]{0.13,0.29,0.53}{\textbf{#1}}}
\newcommand{\DataTypeTok}[1]{\textcolor[rgb]{0.13,0.29,0.53}{#1}}
\newcommand{\DecValTok}[1]{\textcolor[rgb]{0.00,0.00,0.81}{#1}}
\newcommand{\DocumentationTok}[1]{\textcolor[rgb]{0.56,0.35,0.01}{\textbf{\textit{#1}}}}
\newcommand{\ErrorTok}[1]{\textcolor[rgb]{0.64,0.00,0.00}{\textbf{#1}}}
\newcommand{\ExtensionTok}[1]{#1}
\newcommand{\FloatTok}[1]{\textcolor[rgb]{0.00,0.00,0.81}{#1}}
\newcommand{\FunctionTok}[1]{\textcolor[rgb]{0.13,0.29,0.53}{\textbf{#1}}}
\newcommand{\ImportTok}[1]{#1}
\newcommand{\InformationTok}[1]{\textcolor[rgb]{0.56,0.35,0.01}{\textbf{\textit{#1}}}}
\newcommand{\KeywordTok}[1]{\textcolor[rgb]{0.13,0.29,0.53}{\textbf{#1}}}
\newcommand{\NormalTok}[1]{#1}
\newcommand{\OperatorTok}[1]{\textcolor[rgb]{0.81,0.36,0.00}{\textbf{#1}}}
\newcommand{\OtherTok}[1]{\textcolor[rgb]{0.56,0.35,0.01}{#1}}
\newcommand{\PreprocessorTok}[1]{\textcolor[rgb]{0.56,0.35,0.01}{\textit{#1}}}
\newcommand{\RegionMarkerTok}[1]{#1}
\newcommand{\SpecialCharTok}[1]{\textcolor[rgb]{0.81,0.36,0.00}{\textbf{#1}}}
\newcommand{\SpecialStringTok}[1]{\textcolor[rgb]{0.31,0.60,0.02}{#1}}
\newcommand{\StringTok}[1]{\textcolor[rgb]{0.31,0.60,0.02}{#1}}
\newcommand{\VariableTok}[1]{\textcolor[rgb]{0.00,0.00,0.00}{#1}}
\newcommand{\VerbatimStringTok}[1]{\textcolor[rgb]{0.31,0.60,0.02}{#1}}
\newcommand{\WarningTok}[1]{\textcolor[rgb]{0.56,0.35,0.01}{\textbf{\textit{#1}}}}
\usepackage{graphicx}
\makeatletter
\newsavebox\pandoc@box
\newcommand*\pandocbounded[1]{% scales image to fit in text height/width
  \sbox\pandoc@box{#1}%
  \Gscale@div\@tempa{\textheight}{\dimexpr\ht\pandoc@box+\dp\pandoc@box\relax}%
  \Gscale@div\@tempb{\linewidth}{\wd\pandoc@box}%
  \ifdim\@tempb\p@<\@tempa\p@\let\@tempa\@tempb\fi% select the smaller of both
  \ifdim\@tempa\p@<\p@\scalebox{\@tempa}{\usebox\pandoc@box}%
  \else\usebox{\pandoc@box}%
  \fi%
}
% Set default figure placement to htbp
\def\fps@figure{htbp}
\makeatother
\setlength{\emergencystretch}{3em} % prevent overfull lines
\providecommand{\tightlist}{%
  \setlength{\itemsep}{0pt}\setlength{\parskip}{0pt}}
\usepackage{bookmark}
\IfFileExists{xurl.sty}{\usepackage{xurl}}{} % add URL line breaks if available
\urlstyle{same}
\hypersetup{
  pdftitle={Bayes Factors},
  hidelinks,
  pdfcreator={LaTeX via pandoc}}

\title{Bayes Factors}
\author{}
\date{\vspace{-2.5em}2026-04-17}

\begin{document}
\maketitle

\subsection{Introduction}\label{introduction}

A Bayes factor compares two models by the ratio of how plausible the
observed data are under each, after integrating over the parameters of
each model with respect to its prior. It is the Bayesian counterpart of
a likelihood-ratio test, but it carries a richer interpretation: rather
than rejecting or failing to reject a null, it returns a continuous
strength-of-evidence measure that can equally support the null model
when the data are inconsistent with the alternative. This symmetry ---
quantifying support for either side --- is one of the main reasons Bayes
factors recur in applied work where ``no effect'' is itself a
substantively interesting conclusion.

\subsection{Prerequisites}\label{prerequisites}

A working understanding of Bayes' theorem, the role of the marginal
likelihood (evidence) in Bayesian inference, and the difference between
point-null and composite hypotheses.

\subsection{Theory}\label{theory}

For two models \(M_0\) and \(M_1\) with parameter vectors \(\theta_0\)
and \(\theta_1\) and priors \(p(\theta_0 \mid M_0)\) and
\(p(\theta_1 \mid M_1)\), the Bayes factor in favour of \(M_1\) is

\[\mathrm{BF}_{10} = \frac{p(y \mid M_1)}{p(y \mid M_0)} = \frac{\int p(y \mid \theta_1, M_1) p(\theta_1 \mid M_1) \, \mathrm d \theta_1}{\int p(y \mid \theta_0, M_0) p(\theta_0 \mid M_0) \, \mathrm d \theta_0}.\]

Posterior odds equal the Bayes factor times the prior odds; under equal
prior odds, the posterior probability of \(M_1\) is
\(\mathrm{BF}_{10} / (\mathrm{BF}_{10} + 1)\). Jeffreys' interpretive
scale calls \(\mathrm{BF} > 3\) substantial, \(> 10\) strong, \(> 30\)
very strong, and \(> 100\) decisive evidence; the same thresholds apply,
mirrored, for evidence in favour of \(M_0\).

\subsection{Assumptions}\label{assumptions}

Both models are well-specified and proper priors are assigned to all
parameters that differ between the two; improper priors leave the
marginal likelihoods undefined up to an arbitrary constant and the Bayes
factor inherits that indeterminacy. Numerical estimation must be stable
--- bridge sampling, Chib's method, or Savage-Dickey ratios for nested
models.

\subsection{R Implementation}\label{r-implementation}

\begin{Shaded}
\begin{Highlighting}[]
\FunctionTok{library}\NormalTok{(BayesFactor)}

\CommentTok{\# Simple example: one{-}sample t{-}test equivalent}
\NormalTok{x }\OtherTok{\textless{}{-}} \FunctionTok{rnorm}\NormalTok{(}\DecValTok{40}\NormalTok{, }\AttributeTok{mean =} \FloatTok{0.5}\NormalTok{, }\AttributeTok{sd =} \DecValTok{1}\NormalTok{)}
\FunctionTok{ttestBF}\NormalTok{(x)}

\CommentTok{\# Regression}
\FunctionTok{data}\NormalTok{(mtcars)}
\NormalTok{full }\OtherTok{\textless{}{-}} \FunctionTok{lmBF}\NormalTok{(mpg }\SpecialCharTok{\textasciitilde{}}\NormalTok{ wt }\SpecialCharTok{+}\NormalTok{ hp, }\AttributeTok{data =}\NormalTok{ mtcars)}
\NormalTok{reduced }\OtherTok{\textless{}{-}} \FunctionTok{lmBF}\NormalTok{(mpg }\SpecialCharTok{\textasciitilde{}}\NormalTok{ wt, }\AttributeTok{data =}\NormalTok{ mtcars)}
\NormalTok{full }\SpecialCharTok{/}\NormalTok{ reduced}
\end{Highlighting}
\end{Shaded}

\subsection{Output \& Results}\label{output-results}

The \texttt{BayesFactor} package returns a Bayes factor with an
estimated error percentage. For regression-type models the output also
includes the implied posterior model probabilities; ratios of
\texttt{BFBayesFactor} objects deliver pairwise comparisons directly.

\subsection{Interpretation}\label{interpretation}

A reporting sentence: ``A default-prior \(t\)-test Bayes factor was
\(\mathrm{BF}_{10} = 4.7\) (\(\pm 0.001 \%\)), interpreted as
substantial evidence that the population mean differs from zero.'' When
reporting, always include the prior and (for \texttt{BayesFactor}) the
prior scale parameter \texttt{rscale}; readers cannot otherwise judge
how much of the evidence is data-driven versus prior-driven.

\subsection{Practical Tips}\label{practical-tips}

\begin{itemize}
\tightlist
\item
  Wider priors mechanically shift Bayes factors toward the null
  (Lindley--Bartlett paradox); a sensitivity analysis across at least
  two prior scales is mandatory for any Bayes-factor conclusion.
\item
  \texttt{bridgesampling::bridge\_sampler()} works directly with
  \texttt{brms} and Stan fits; it is the most general way to compute
  marginal likelihoods for non-trivial models.
\item
  For nested null hypotheses, the Savage-Dickey density ratio is faster,
  more stable, and avoids marginal-likelihood estimation entirely.
\item
  Bayes factors near 1 are genuinely inconclusive; resist the temptation
  to declare evidence with \(\mathrm{BF} = 1.5\) in either direction.
\item
  Bayes factors are not a substitute for predictive performance; pair
  them with \texttt{loo()} when comparing models for prediction, since
  the two answer different questions.
\end{itemize}

\subsection{R Packages Used}\label{r-packages-used}

\texttt{BayesFactor} for default-prior tests and ANOVA / regression
Bayes factors, \texttt{bridgesampling} for general marginal-likelihood
estimation, \texttt{bayestestR} for tidy summaries that work with
\texttt{brms} and \texttt{rstanarm} fits, and \texttt{polspline} for
Savage-Dickey density estimation when the null is nested.

\subsection{For Reviewers}\label{for-reviewers}

\textbf{What to look for in a paper using this method.}

\begin{itemize}
\tightlist
\item
  \textbf{Common misapplications.}
\item
  \textbf{Diagnostics that should be reported but often aren't.}
\item
  \textbf{Red flags in tables and figures.}
\item
  \textbf{What to verify.}
\item
  \textbf{What an adequate Methods paragraph must contain.}
\end{itemize}

\subsection{See also --- labs in this
chapter}\label{see-also-labs-in-this-chapter}

\begin{itemize}
\tightlist
\item
  \href{labs/lab-bayesian-thinking-control-vs-decision-error.Rmd}{Bayesian
  thinking; α control vs decision error}
\item
  \href{labs/lab-brms-stan-loo-hierarchical-models.Rmd}{brms / Stan,
  LOO, hierarchical models}
\end{itemize}

\end{document}
